\section{A Benchmark Agenda}


A taxonomy earns its keep by becoming a suite. We propose one scene family per type (Fig. \ref{fig:pipeline}), each shipping three things:

\begin{enumerate}
  \item \textbf{A blind-policy defect measurement} --- the success-conditioned violation rate of an unmodified VLA, with confidence intervals.
  \item \textbf{Fixability ablations} --- name the hazard vs not, render it vs hide it, which locate the failure on the prompting / perception / architecture spectrum.
  \item \textbf{A reference safety layer} --- an external baseline (repulsion shield, protective stop, speed governor) with its efficacy \emph{and failure modes} reported.
\end{enumerate}

Two further items belong in every type's protocol: a \textbf{feasibility witness} --- a trajectory that completes the scene without violating the predicate, so that a violation rate is attributable to the policy rather than the scene --- and a \textbf{protective-stop reference layer}, the reactive response the standards prescribe, reported next to the repulsion shield --- run here for T6 and T3a, where the stopped T6 carry and the shielded-and-governed T3a carry are the first two witnesses (\S5.4, \S5.7). \textbf{Release.} Scenes, recorders, analysis scripts and every per-episode log behind Appendix A are released (anonymized repository); adding a policy is a server swap.

Canonical thresholds and the open design decisions are in Appendix H.

